\scanpage{8}
\mspage{6}

morphismes de flêches.

Notons une symmétrie évidente dans la \unsure{catégorie|c atégorie} (Bicat), interprété
en termes d\deleted{\unsure{a|e}}'ensemble bi-semisimpliciaux, obtenue en échangeant les lign
et les colonnes dqns \deleted{les} de tels ensembles bi... \unsure{\deleted{q}U|\deleted{U}q}n ensemble bisemisimpl
al qui ... peut donc de deux façons différentes être consi\deleted{i}déré comme déf
nissant une bicatégorie. De même le groupe symétrique $\mathrm{Sym}_n$ opère sur
la catégorie \deleted{de} (\unsure{nupcat|n-upcat}) des catégories n-uples ..

Une 2-catégorie peut être définie comme constituée par un ensemble
\deleted{discret} $C^{\deleted{\unsure{o}}}_0$ (l'ensemble des objets ), qu'on interprètera comme une catégo%
rie discrète, plus la donnée pour tout $x,y \in \unsure{C^o_0}$ d'une catégorie $\underline{\mathrm{Hom}}(x,y)$,
plus \deleted{une} \inserted{des foncteurs} composition entre celles\unsure{\deleted{(}-|\deleted{-}(}ci ... Prenant la catégorie somme \deleted{de} C
des $\underline{\mathrm{Hom}}(x,y)$, On peut exprimer la donnée de ces foncteurs de composition
et des sources et buts\begin{tabular}[b]{@{}l@{}}\inserted{$x,y\unsure{)}$}\\[-.6ex]\begin{tikzpicture}[x=1mm,y=1mm,baseline=0pt]\draw[line width=.5pt] (0,0)..controls(2,1.6)and(4.7,3.2)..(6.8,4.2);\draw[line width=.3pt] (7.1,4.4)--(7.5,4.7);\draw[line width=.3pt] (7.7,4.9)--(8.4,5.7);\draw[line width=.3pt] (8.6,5.9)--(9,6.5);\end{tikzpicture}\end{tabular} dans les catégories $\underline{\unsure{\deleted{P}H|\deleted{R}H|\deleted{H}P|\deleted{H}R}\mathrm{om}}(x,y)$, comme définis par
des foncteurs ``source'' et ``but'' $C^{\deleted{1}}_1 \to \unsure{C^o_0}$, et un foncteur de composition
$C^{\deleted{1}}_1 \mathrm{x}_{C^o_0} C^{\deleted{1}}_1 \longrightarrow \unsure{C^{\deleted{1}}_1}$ , les axiomes des 2-catégories correspondant exactement au
axiomes des(Cat)-catégories. De cette façon, la catégorie des 2-catégor
est équivalente à la catégorie des (Cat)-catégories qui ont une ``catégo
d'objets'' qui est discrète, ou encore à la catégorie des Bicatégories
qui, interprétées comme ensembles bi-semisimpliciaux, ont \deleted{une} comme
première ligne un ensemble bemis\unsure{\deleted{s}i|\deleted{i}s}\unsure{\deleted{i}m|\deleted{m}i}plicial discret (i.e. un foncteur cons%
tant sur $\Delta$ ). Il parait qu'on a une interprétation simple dans le même
esprit de la notion de n-catégorie en termes d'ensembles n-semi-simplici
\margin{\begin{tikzpicture}[x=1mm,y=1mm]\draw[line width=.45pt] (0,21).. controls (-0.4,15) and (0.7,8)..(0.7,3).. controls (0.7,1.4) and (0.9,1.2)..(1.1,1.2); \draw[line width=.45pt] (1.2,20.4).. controls (1.3,14) and (2.1,7)..(2.5,3).. controls (2.6,1.5) and (3.4,0.2)..(3.2,0).. controls (3.1,-0.1) and (2.9,0)..(2.9,0.1); \end{tikzpicture}}
\deleted{\unsure{ou|au}} Demander à Quillen relations entre la notion de m-Catégorie dans
la catégorie (n-Cat), et celle d\deleted{\illegible{1}}\inserted{'} ensemble (m+n)-semisimplicial.
\deleted{(en particulier entre cette dernière notion et celle de}

L'interprétation standard semisimpliciale permet d'interpréter

\margin{7}
